\documentclass[11pt]{article}

\usepackage[a4paper,margin=2.4cm]{geometry}
\usepackage{lmodern}
\usepackage[T1]{fontenc}
\usepackage[utf8]{inputenc}
\usepackage{microtype}
\usepackage{graphicx}
\usepackage{xcolor}
\usepackage{float}
\usepackage{caption}
\usepackage{enumitem}
\usepackage{parskip}
\usepackage{titlesec}
\usepackage[hidelinks]{hyperref}

% Same screenshots as the English guide and the in-app help page.
\graphicspath{{../../frontend/public/help/}}

\definecolor{accent}{HTML}{1D4ED8}
\definecolor{accentbg}{HTML}{EFF4FF}
\definecolor{soft}{HTML}{555555}

\titleformat{\section}{\Large\bfseries\color{accent}}{}{0pt}{}
\titlespacing*{\section}{0pt}{1.4em}{0.5em}

\captionsetup{font={small,it},labelformat=empty,skip=4pt}
\setlength{\fboxsep}{10pt}

\newcommand{\shot}[2]{%
  \begin{figure}[H]\centering
  \includegraphics[width=\linewidth,height=0.6\textheight,keepaspectratio]{#1}%
  \caption*{#2}%
  \end{figure}}

\newcommand{\callout}[1]{%
  \begin{center}
  \fcolorbox{accent}{accentbg}{\parbox{0.96\linewidth}{\small #1}}
  \end{center}}

\hypersetup{
  pdftitle={Minerva: Lärarguide},
  pdfauthor={DSV, Stockholms universitet},
}

\begin{document}

\begin{center}
  {\LARGE\bfseries Minerva}\\[2pt]
  {\Large\color{accent} Lärarguide}\\[6pt]
  {\color{soft}\small Koppla AI-studieassistenten Minerva till din Moodle-kurs \textbullet{} ungefär 5 minuter \textbullet{} 3 steg}
\end{center}

\vspace{0.4em}
\hrule height 1pt
\vspace{1.0em}

\noindent
Minerva är en AI-studieassistent som svarar utifrån ditt eget kursmaterial. Din
kurs, dess material och dess inspelade föreläsningar hålls redan uppdaterade i
Minerva, så det som återstår är att lägga in en länk till den i Moodle.

\callout{\textbf{Innan du börjar:} din administratör har registrerat Minerva som
ett externt verktyg i hela Moodle. Du behöver fortfarande aktivera det för var
och en av dina kurser (Steg 1) innan det visas i aktivitetsväljaren. Om du inte
ser alternativen som visas här, kontakta
\href{mailto:lambda@dsv.su.se}{lambda@dsv.su.se}.}

\section{Det här sker automatiskt}

\begin{itemize}[leftmargin=1.4em,itemsep=4pt]
  \item Kurser importeras från Daisy varje natt, så din kurs finns redan i
  Minerva.
  \item En Moodle-kurs vars \textbf{Course ID number} innehåller Daisy-tillfällets
  id kopplas automatiskt till motsvarande Minerva-kurs, oftast inom en timme.
  \item Kopplat material synkas två gånger i timmen, så nya filer och sidor kommer
  med utan att du gör något.
  \item Inspelade föreläsningar från DSV Play hämtas via kursens Play-beteckning,
  transkript och allt.
\end{itemize}

\shot{minerva-dashboard.png}{Din Minerva-översikt med de kurser du undervisar i.}

\section{Steg 1. Aktivera Minerva i din kurs}

Minerva är registrerat i hela Moodle, men det visas inte i en kurs förrän du
aktiverar det. I din Moodle-kurs, öppna menyn \textbf{More} och välj \textbf{LTI
External tools}. Hitta \textbf{Minerva}, aktivera det och slå på \textbf{Show in
activity chooser}.

\shot{enable-tool.png}{Kursens LTI External tools: aktivera Minerva och visa det i aktivitetsväljaren.}

När det är korrekt aktiverat visas Minerva i aktivitetsväljaren med sin logotyp.
Om du istället ser en tom eller generisk ikon har verktyget registrerats utan sin
ikon; kontakta \href{mailto:lambda@dsv.su.se}{lambda@dsv.su.se}.

\section{Steg 2. Lägg till Minerva-aktiviteten i ett avsnitt}

I din Moodle-kurs, slå på \textbf{Edit mode}, välj \textbf{Add an activity or
resource} i det avsnitt där du vill ha assistenten och välj \textbf{Minerva}.

\shot{activity-chooser.png}{Välj Minerva i aktivitetsväljaren; den har Minervas logotyp.}

Ge den ett namn som studenterna känner igen (till exempel ``Minerva AI
assistant'') och välj sedan \textbf{Save and return to course}. Aktiviteten visas
nu i din kurs.

\shot{lti-config.png}{Namnge aktiviteten och spara.}
\shot{course-activity.png}{Minerva-aktiviteten finns nu i din kurs.}

\section{Steg 3. Öppna den en gång och välj kurs}

Klicka på Minerva-aktiviteten du lade till. Första gången väljer du vilken
Minerva-kurs den ska öppna och väljer \textbf{Link and open Minerva}. Minerva
kommer ihåg detta, så varje senare öppning av dig eller dina studenter går till
rätt kurs.

\shot{first-launch-bind.png}{Vid första öppningen kopplar du aktiviteten till din Minerva-kurs.}

\section{Klart: så här ser studenterna det}

Studenterna klickar på samma aktivitet och får en chattassistent som svarar
utifrån ditt kursmaterial, direkt i Moodle. De föreslagna frågorna hämtas från
materialet.

\shot{student-chat.png}{Studentvyn: en kursförankrad chattassistent.}
\shot{launch-in-moodle.png}{Samma assistent, inbäddad i Moodle-aktiviteten.}

\section{Bestäm exakt vad som indexeras}

Den automatiska kopplingen indexerar det som finns i din Moodle-kurs. Vill du
istället välja precis vad assistenten svarar utifrån skapar du en egen kurs i
Minerva med \textbf{New Course} och laddar upp bara de dokument du vill ha med.

Lägg sedan till ytterligare en Minerva-aktivitet i Moodle och peka den mot den
kursen första gången du öppnar den. Båda aktiviteterna kan ligga i samma
Moodle-kurs; studenterna får den du skickar dem till.

\section{Om den automatiska kopplingen uteblev}

Utan ett Daisy-tillfälles id i Moodle-kursens \textbf{Course ID number} finns
inget att matcha mot, och då kopplas ingenting. Öppna \textbf{Minerva settings}
via kursens \textbf{More}-meny, välj Minerva-kursen och klicka \textbf{Link
Minerva course}.

\shot{link-form.png}{Minerva settings: välj din Minerva-kurs och länka den.}

\textbf{Sync materials} laddar då upp dina filer, sidor och
avsnittssammanfattningar direkt; därefter håller synken två gånger i timmen dem
aktuella.

\shot{linked.png}{När kursen är länkad kan du synka material eller avlänka när som helst.}
\shot{sync-complete.png}{Minerva bekräftar hur många resurser som laddades upp.}

\section{Bra att veta}

\begin{itemize}[leftmargin=1.4em,itemsep=4pt]
  \item Assistenten svarar bara utifrån materialet i den Minerva-kurs som
  aktiviteten pekar på. Studenternas konversationer är synliga för kurspersonal;
  se \textbf{Data handling} i Minerva för detaljer.
  \item En kurs som kopplats automatiskt visar det i \textbf{Minerva settings}:
  tar du bort länken där förblir den olänkad och kopplas inte automatiskt igen.
  \item Lärarbesvarade forumtrådar kan tas med från \textbf{Minerva settings}
  (valfritt; studenternas namn tas bort före uppladdning).
  \item Frågor eller behöver du högre användningsgränser? Mejla
  \href{mailto:lambda@dsv.su.se}{lambda@dsv.su.se}.
\end{itemize}

\vspace{1.2em}
\hrule height 0.4pt
\vspace{0.4em}
\noindent{\color{soft}\small Minerva \textbullet{} Institutionen för data- och
systemvetenskap (DSV), Stockholms universitet \textbullet{}
\href{mailto:lambda@dsv.su.se}{lambda@dsv.su.se}}

\end{document}
